\clearpage
\fancyhead[L]{\small IRIS.}
\SourceLeaf{480}
\begin{center}
\rule{0.55\linewidth}{0.4pt}\\[2pt]
Astronomische Nachrichten. Band XXVI. Nr. 616. Seite 241. u. 255. 1847 November 8.\\[2pt]
\rule{0.55\linewidth}{0.4pt}
\end{center}
\begin{center}[\textsc{Gauss} an \textsc{Schumacher}.]\quad Göttingen 1847 October 11.\end{center}

Ich kann Ihnen heute doch noch eine Meridianbeobachtung der \emph{Iris} mittheilen, nach einer Unterbrechung von vollen drei Wochen. Der Planet war aber wegen seiner Lichtschwäche sehr schwer zu beobachten.

\begin{center}
\begin{tabular}{ll}
\multicolumn{2}{l}{1847 Oct. 10.\quad $6^{\rm u}\ 42^{\rm m}\ 2^{\rm s}6$ M. Z.} \\
Ger. Aufst. & $299^\circ\ 13'\ 32''7$ \\
Abweichung südl. & $14^\circ\ 19'\ 38''6$ \\
\end{tabular}
\end{center}

\textsc{Goldschmidts} Elemente stimmen noch sehr gut. Von der \emph{Sonnenfinsterniss} habe ich nur das Ende beobachten können.

\begin{center}
1847 Oct. 8.\quad $21^{\rm u}\ 39^{\rm m}\ 33^{\rm s}7$ Mittl. Zeit.\\
Herr Prof. \textsc{Goldschmidt} $2^{\rm s}4$ später.
\end{center}

Ausserdem habe ich noch den Eintritt eines Fleckens, d. i. Anfang des Eintritts, oder erste Berührung aufgezeichnet $19^{\rm u}\ 3^{\rm m}\ 40^{\rm s}8$, aber bei so stark wallender Luft, dass das deutliche Sehen sehr dadurch beeinträchtigt wurde. Prof. \textsc{Goldschmidt} hatte $3^{\rm s}7$ weniger.

In einem späteren Schreiben fügte der Herr Geheime Hofrath noch folgende Meridianbeobachtungen der \emph{Iris} hinzu:

\begin{center}
\setlength{\tabcolsep}{4pt}
\begin{tabular}{l|r|r|r}
1847 & Mittlere Zeit & Gerade Aufsteig. & Südl. Abweich. \\\hline
October 11 & $6^{\rm u}\ 39^{\rm m}\ 0^{\rm s}6$ & $299^\circ\ 27'\ 3''6$ & $14^\circ\ 18'\ 39''3$ \\
\hfill 13 & $6\ 33\ 1.5$ & $299\ 55\ 18.2$ & $14\ 16\ 11.0$ \\
\hfill 14 & $6\ 30\ 4.2$ & $300\ 9\ 59.4$ & $14\ 14\ 44.5$ \\
\end{tabular}
\end{center}

\clearpage
\hypertarget{native-480-october21}{}
\begin{center}[\textsc{Gauss} an \textsc{Schumacher}.]\quad Göttingen 1847 October 21.\end{center}

Ich kann Ihnen noch drei Meridianbeobachtungen der \emph{Iris} schicken; ich schreibe jedoch auch die Beobachtungen vom 13ten und 14ten noch einmal hieher, da die Declinationen in Folge der neuen Nadirbestimmungen eine geringe Aenderung erlitten haben, und ich nicht weiss, ob dieselbe an den Ihnen früher mitgetheilten Zahlen schon angebracht war. Die Angaben von 10 und 11 werden davon nicht afficirt.

\begin{center}
\setlength{\tabcolsep}{4pt}
\begin{tabular}{l|r|r|r}
1847 & Mittlere Zeit & Gerade Aufsteig. & Südl. Abweich. \\\hline
October 13 & $6^{\rm u}\ 33^{\rm m}\ 1^{\rm s}5$ & $299^\circ\ 55'\ 18''1$ & $14^\circ\ 16'\ 11''4$ \\
\hfill 14 & $6\ 30\ 4.2$ & $300\ 9\ 59.4$ & $14\ 14\ 45.1$ \\
\hfill 15 & $6\ 27\ 8.2$ & $300\ 25\ 1.5$ & $14\ 13\ 17.2$ \\
\hfill 16 & $6\ 24\ 13.9$ & $300\ 40\ 27.1$ & $14\ 11\ 50.3$ \\
\hfill 17 & $6\ 21\ 20.8$ & $300\ 56\ 12.1$ & $14\ 10\ 9.1$ \\
\end{tabular}
\end{center}

Seitdem ist es um die Culminationszeit immer trübe gewesen, was ich um so mehr bedauert habe da ich dadurch verhindert bin, den Planeten in sehr naher Zusammenkunft mit einem Sterne 9ter Grösse aus \textsc{Bessels} Zonen zu sehen (20, 101 \textsc{Weisse}, dem der Planet am 18ten fast auf eine halbe Minute nahe gekommen sein muss). Uebrigens war der Planet wegen seiner grossen Lichtschwäche immer sehr schwer zu beobachten, ich kann aber nicht sagen, dass die grosse Nähe des Mondes am 17ten die Schwierigkeit des Beobachtens grösser gemacht hat, als sie an den frühern Abenden war.
